% RESEARCH_DRAFT_SHA256: 066526c6fa3f758f17e6d5e0a5c7ea2e7c18e1fc1b7e3de4c80957c2fb8c1066
% Layout and source-reference wrappers expanded for the web renderer.
\documentclass[11pt]{article}
\usepackage[margin=25mm]{geometry}
\usepackage{fontspec,unicode-math}
\setmainfont{Latin Modern Roman}
\setmathfont{Latin Modern Math}
\usepackage{amsmath,amssymb,mathtools,xurl,hyperref}
\hypersetup{colorlinks=true,urlcolor=blue,linkcolor=blue}
\setcounter{secnumdepth}{0}
\setlength{\parindent}{0pt}
\setlength{\parskip}{5pt}
\setlength{\emergencystretch}{3em}
\newcommand{\sref}[2]{\hyperlink{source:#1:#2}{[S#2]}}
\newcommand{\cataloguescope}[1]{\paragraph{Source record.}#1}
\newcommand{\floor}[1]{\left\lfloor#1\right\rfloor}
\begin{document}
\section*{Research-attempt provenance}
This standalone research attempt retains the complete original definition body below,
including its historical source and preparation notes. Those notes describe the
earlier catalogue preparation; the new attempt and its verification are identified
separately after that body. No original ID, catalogue status, or source statement is changed.

% UnsolvedMath ID 20002172; source catalogue code AIM-INFRASTRUCTURE-0072
\section{Profile-scoped evidence for PreTeXt interactivity and accessibility}\label{20002172}
{\small UnsolvedMath ID 20002172 · Computer Science}
\subsection{Definitions and mathematical statement}

\paragraph{Shared notation.}Standard number systems are $\mathbb N,\mathbb Z,\mathbb Q,\mathbb R,\mathbb C$. Specialized notation and index ranges follow the cited source and the definitions below.


Interactivity lets a reader change inputs or receive responses while using a document. Accessibility concerns access to content and controls across differing user needs and technologies. Semantic document structure can support multiple presentations, but a format alone does not establish the usability of every generated document.
\paragraph{Record scope and qualifications.}
The catalogue title is retained as metadata. Named certificates, contracts, bounds or frameworks in that title are not asserted as results of the original source extract.
The source is an authoring-advocacy question. It supplies no benchmark, audience, or universal accessibility guarantee.

Catalogue screening limitation: PreTeXt software design, documentation, or publication feature request; no mathematical problem is specified by the supplied passage.

The source snapshot labels this record partially solved; this is catalogue metadata rather than independent verification of a mathematical result or the current completion of a project.
\paragraph{Statement recorded in the supplied source.}
Why write in PreTeXt? Interactivity, Accessibility. \href{https://github.com/PreTeXtBook/community/wiki/Moderated-Problem-Session}{[S2]}
\paragraph{Source status.}
The supplied catalogue labels this source record partially solved; that metadata was checked against the supplied snapshot on 2026-10-11. The label is preserved for provenance and is not an independent verification of present open status or mathematical correctness.
\subsection{Short English statement}
Explain the benefits and requirements of PreTeXt authoring, including interactivity and accessibility.
\subsection{Sources}
\begin{itemize}
\item[S1] Ulam AI and the UnsolvedMath Contributors, UnsolvedMath: A Curated Collection of Open Mathematics Problems, record 20002172 (AIM-INFRASTRUCTURE-0072). Catalogue curation and metadata: CC BY 4.0. The source extract is retained; definitions, English summaries and qualifications are editorial additions.
\url{https://huggingface.co/datasets/ulamai/UnsolvedMath}.
\item[S2] AIM Workshop Problem Lists, source linked by catalogue record AIM-INFRASTRUCTURE-0072.
\url{https://github.com/PreTeXtBook/community/wiki/Moderated-Problem-Session}.
\end{itemize}
\paragraph{Preparation and review.}
The original source and UnsolvedMath contributors retain attribution for the source extract. Codex (GPT-6; exact serving revision not exposed) assisted with definitions, English summaries, scope notes and formatting. No new human mathematical authorship or independently verified solution is asserted. Permanent source ID 20002172 and code AIM-INFRASTRUCTURE-0072 checked against the complete supplied registry and current repository. Distinct upstream records, including repeated questions and broad or damaged extracts, are retained for complete coverage of the source CS category. A repeated record is not a new independent conjecture. Complete mathematical-equivalence and current-literature review has not been performed. Preparation and final review follow contribution prompts C and D, with source gaps disclosed. Structural and rendering checks do not establish a new theorem or open status.
\label{end:20002172}

\section{Literature and context check}
The primary community-session source was consulted. The record is a broad software-design request, and no mathematical theorem follows from its editorial catalogue title. The finite model below is explicitly auxiliary; it does not certify an actual implementation or current feature status.
\section{Research attempt: A visual control can fail a keyboard reachability test}
Select one narrow accessibility-related task: every essential control must be reachable and activatable from the initial state using keyboard actions. Model the interface as a finite directed graph whose edges are labeled keyboard or pointer. Suppose the initial state $s$ has a pointer-only edge to a submit state $t$, and no keyboard edge reaches $t$. The control can be visible, correctly named, and mouse-operable while failing the selected keyboard task. Thus visual presence and text alternatives alone do not prove this task property.

For the finite graph, delete all non-keyboard edges and compute the reachable set $R$ from $s$. A required target $t$ is keyboard-reachable exactly when $t\in R$: graph-search insertion produces an actual path to every inserted state, and induction on path length proves that every keyboard path endpoint is inserted. Retaining predecessor edges gives a witness sequence when the check succeeds, or the absence of any reachable target when it fails. If activation itself is modeled as a required edge, the check must include that edge's action and resulting state, not merely focus on the control.

This property is only one component of an accessibility evaluation. Add a second requirement that a result be announced to a screen reader. Two systems can have identical keyboard-transition graphs but different announcement outputs; one emits a result message and the other emits silence. The reachability test cannot distinguish them because its input omitted announcements. To verify both tasks, the model must contain both transition and observation data, and the profile must say which observations count. A statement about a markup element or a finite test profile therefore cannot certify every user's interaction experience.
\section{Verification and adversarial audit}
The construction and counterexample are checked in the written finite model. No actual PreTeXt build, browser, EPUB reader, or publishing service was tested. All runtime and identity assumptions used by the proof are stated explicitly.
\paragraph{Audit corrections.} A narrow keyboard criterion is proved without being relabeled full accessibility. Announcement behavior supplies a concrete counterexample to that extrapolation.
\paragraph{External computation.} None; symbolic arguments and boundary checks only.
\section{Scope and remaining gaps}
This is a restricted route investigation. The full feature requires implementation, supported-version decisions, conformance tests, and user-facing evaluation that were not performed here. The original agenda remains unresolved, and no novelty or completed software feature is claimed.
\paragraph{Integration review.} Codex (GPT-6; exact serving revision not exposed) reviewed the written arguments and expanded layout and source-reference wrappers. Repository metadata and the production web renderer were checked using Python 3.13 and JavaScript with MathJax 3.2.2. This model-assisted review is not independent human proof verification; no PDF compilation or formal proof checking is asserted.

\section{Completion Estimate}
\noindent\textbf{COMPLETION: 0\%}
The calculations below the preserved source statement investigate a concrete route,
but establish no substantiated new progress toward the entire original target.
Known special cases and failed sufficient conditions are not credited as a solution.
\section{References}
The original source references and catalogue attribution are preserved in the
Sources section above. The following URLs were actually consulted for this attempt;
retention of an additional original reference does not claim fresh consultation of it.
\begin{itemize}
\item \url{https://github.com/PreTeXtBook/community/wiki/Moderated-Problem-Session} (consulted 2026-10-11). 
\end{itemize}
\section{Model and contributor disclosure}
OpenAI gpt-6-astra, requested reasoning configuration ultra, generated this research
attempt on 2026-10-11. The exact serving revision is not exposed. Codex GPT-6
performed orchestration and repository-format review; its exact serving revision
is also not exposed. The mathematical route, boundary checks and final adversarial
self-audit followed contribution prompts A then B in
\texttt{docs/contribution\_prompts.md}, after reading \texttt{CONTRIBUTING.md}.
Python 3.14.2 assembled the standalone documents; any mathematical computation is
specified in Verification. No theorem prover was used. No human mathematical
authorship, substantial human proof correction, or independent human proof checking
is asserted. Model-assisted review, source retrieval, and structural checks do not
establish correctness or novelty. Full private reasoning transcripts are unavailable
as reproducibility artifacts; the complete arguments, inputs and outcomes actually
used are stated here or in the companion batch audit.
\section{Final status}
Unresolved as a new research contribution. This is a concrete route investigation,
not a claim of new complete resolution or a current-status certification. A known
answer reproduced for context does not become a new contribution. The catalogue's
source-status label is preserved; this attempt status does not assert that a
source-labelled solved question is open.
\end{document}
